\section{Problem Setup and Evaluation Targets}
\label{sec:setup}
\subsection{Prescribed actions and fixed representations}

\paragraph{Prescribed actions and paired states.}
Let $S$ be a finite set of $n$ states and $G$ a finite group.
The \emph{prescribed action} is the homomorphism $\alpha:G\to\mathrm{Sym}(S)$ that specifies the intended transformations of the named states.
Write $g(i)=\alpha(g)(i)$. Then $e(i)=i$ and $g_1(g_2(i))=(g_1g_2)(i)$ for every $i\in S$ and $g_1,g_2\in G$ \citep{serre1977linear}.
For the cyclic group $C_7$, the generator advances the day by one, and the action on the seven days is free and transitive.
The eight rotations and reflections in $D_4$ act faithfully and transitively on the four corners NW, NE, SE, and SW.
The identity and one diagonal reflection fix each corner, so the $D_4$ action is not free.

Let $x_{t,i}$ denote the prompt obtained by inserting state $i$ into template $t\in\mathcal T_{\mathrm{tpl}}$.
The input action in the Introduction is $g\cdot x_{t,i}=x_{t,g(i)}$, which changes the state while holding the template fixed.
Each template supplies all $n$ state variants.
The paired design includes fixed points, for which the source and target coincide. Each evaluation below specifies which pairs enter its loss.
Study~1 uses 40 templates per domain, split into 28 training, six validation, and six test templates in each partition.
Study~2 uses 48 training, 16 validation, and 64 test context families, each with two template variants.
\Cref{sec:design} gives the partition details.

\paragraph{State-span readout and fixed coordinates.}
For a given model and layer, $h_{i,t}=f(x_{t,i})\in\R^{1\times d}$ is the mean of the layer's block outputs at the token positions whose character offsets overlap the named state. The span can cover one or more tokens.
Causal attention restricts this readout to the tokenized prefix through the last selected token. When that prefix is unchanged, later tokens do not affect the readout.
The splits hold out templates or context families. They do not enforce disjoint tokenized prefixes.
The equivariance target $f(g\cdot x_{t,i})$ is therefore $h_{g(i),t}$.
We store readouts as row vectors so that fitted matrices act by right multiplication.

For each domain, model, layer, and split, we use the complete training set of state--template pairs to estimate the mean and fit principal component analysis (PCA):
\[
 \hat\mu=\frac{1}{n|\mathcal T_{\mathrm{train}}|}
 \sum_{t\in\mathcal T_{\mathrm{train}}}\sum_{i\in S}h_{i,t}
 \in\R^{1\times d},
\]
where $\mathcal T_{\mathrm{train}}\subset\mathcal T_{\mathrm{tpl}}$ is the set of training templates. The projected readout is
\begin{equation}
 z_{i,t}=(h_{i,t}-\hat\mu)U_{\mathrm{PCA}}\in\R^{1\times q},
 \qquad U_{\mathrm{PCA}}^\top U_{\mathrm{PCA}}=I_q.
 \label{eq:projection}
\end{equation}
The columns of $U_{\mathrm{PCA}}\in\R^{d\times q}$ are the leading eigenvectors of the centered training covariance.
We use the training mean and basis on validation and test templates as well.
The fitted predictors in the next subsection act on $z_{i,t}$.
The class-risk analysis applies a separate second-moment whitening step.

PCA captures total training variation, including variation across template means and variation between states within a template.
For weekdays, the seven state means averaged over templates span at most six dimensions after centering.
Within any one template, the centered state differences also span at most six dimensions, but their directions can vary across templates.
The fraction of total variance retained therefore does not tell us how much within-template state variation the projection preserves.

We additionally assess sensitivity to within-template centering
in the saved 16-dimensional Study~2 coordinates.
On each complete test bank, we subtract the mean across states
within each template and compare retaining the original
test-bank whitening with using the centered bank's own whitening
(\cref{app:template-centering-sensitivity}).
This comparison characterizes state-dependent variation
within the retained PCA subspace.
Measuring the fraction of full-dimensional within-template
variation preserved by that subspace, or fitting PCA specifically
to such variation, requires the corresponding full-dimensional
readouts.

Study~1 fixed $q=16$ before the final structural run.
Study~2 retained this dimension for comparability with its development analyses and fitted PCA on its own training templates.
The Study~2 spectral sensitivity analysis uses the first eight saved coordinates to assess dependence on the retained dimension (\cref{app:spectral-dimension-sensitivity}).
The lexical comparison matches the training projected ranks at $q=6$ for weekdays and $q=3$ for corners.
\Cref{app:dimension-rank-audit} documents the choice of dimension, the historical dimension sweep, and the evidence for the rank match.

\subsection{What a fitted operator predicts}

This subsection describes fitted-operator evaluation in the
structural recovery study (Study~1).
For each domain, model, layer, and training split, we use the
training-PCA coordinates $z_{i,t}$ defined in
\cref{eq:projection}.
A fitted matrix $R_g\in\R^{q\times q}$ predicts the transformed
coordinate by right multiplication,
\begin{equation}
 z_{i,t}R_g\approx z_{g(i),t}.
 \label{eq:prediction}
\end{equation}
The maps have no intercept and act in the PCA coordinates
before any second-moment whitening.
Within each condition, the same matrix for an element is
applied across all templates.
Evaluation holds the training mean, PCA basis, and fitted
matrices fixed and uses complete held-out templates containing
the same state vocabulary.

Independent ridge fits a separate matrix for every group
element, including the identity.
Joint fitting optimizes all element matrices together,
combining a prediction loss with soft penalties for the
group law, generator relations, identity, and orthogonality.
Its identity matrix is learned subject to the identity penalty.
The fitting objectives and their weights are specified in
\cref{app:methods}; the normalized evaluation errors are
defined below.
Under the row-vector convention, applying $g_2$ and then $g_1$
multiplies a source row first by $R_{g_2}$ and then by $R_{g_1}$.
An exact row action therefore obeys
\begin{equation}
 R_{g_2}R_{g_1}=R_{g_1\circ g_2},
 \qquad
 (g_1\circ g_2)(i)=g_1(g_2(i)).
 \label{eq:convention}
\end{equation}
We evaluate departures from this law both on the observed
coordinates and in the ambient matrix space.

The generator-only comparison fits the one-day shift $r$
for $C_7$, or the clockwise quarter-turn $r$ and the
vertical-axis reflection $s$ for $D_4$.
The latter swaps NW with NE and SW with SE.
For a fixed word $g=g_1\circ\cdots\circ g_k$, the predicted
target is $z_{i,t}R_{g_k}\cdots R_{g_1}$.
An exact action assigns the same matrix to every word
representing the same element, whereas fitted generator
products can depend on the chosen word.
The weekday evaluation uses the five excluded elements
$r^2,\ldots,r^6$, with $r^7$ evaluated separately as a
cycle diagnostic.
The corner evaluation uses a fixed shortest generator word
for each of the five elements in $G\setminus\{e,r,s\}$.
These products are scored against the held-out target
coordinates $z_{g(i),t}$.
This comparison holds out both the fitted action elements
and the evaluation templates while retaining the state
vocabulary.
The word construction and matched comparators are described
in \cref{app:methods}.

For paired row matrices $A,B\in\R^{N\times q}$, Study~1 uses
\begin{equation}
 D(A,B)=\frac1N\sum_{j=1}^{N}
 \frac{\|A_j-B_j\|_2^2}{\|B_j\|_2^2+10^{-12}}.
 \label{eq:D}
\end{equation}
This averages squared discrepancies after normalizing each
row by the energy of its comparison target.
Let $Z_t$ stack all $n$ source-state rows in test template $t$,
and let $Z_{g,t}$ stack their corresponding transformed targets
in the same source-state order.
For the set $\mathcal T_{\mathrm{test}}$ of held-out templates,
the prediction and data-space consistency errors are
\begin{align*}
 \Eeq
 &=
 \frac{1}{|\mathcal T_{\mathrm{test}}|(|G|-1)}
 \sum_{t\in\mathcal T_{\mathrm{test}}}
 \sum_{\substack{g\in G\\g\ne e}}
 D(Z_tR_g,Z_{g,t}),\\
 \Ehom
 &=
 \frac{1}{|\mathcal T_{\mathrm{test}}||G|^2}
 \sum_{t\in\mathcal T_{\mathrm{test}}}
 \sum_{g_1,g_2\in G}
 D\!\left(
   Z_tR_{g_2}R_{g_1},
   Z_tR_{g_1\circ g_2}
 \right).
\end{align*}
The prediction average excludes the identity element and
retains fixed-point pairs of nonidentity elements.
In $\Ehom$, the comparison target and its normalization
come from the directly fitted prediction
$Z_tR_{g_1\circ g_2}$.
The metric $\Ecomp$ uses the same consistency expression,
averaged over $g_1,g_2\in G\setminus\{e\}$;
their product may still equal the identity.
These averages are computed within each model--layer--split
condition before the higher-level aggregation described in
\cref{app:methods}.
The relationship between prediction and consistency also
depends on operator norms and the comparison denominators
(\cref{prop:propagation}).
The operator-space homomorphism diagnostic uses an unsquared
relative Frobenius norm.
\Cref{app:metrics} specifies the relation diagnostics,
including their comparison targets and aggregation rules
in data and operator space.
The within-template centering analysis in
\cref{app:template-centering-sensitivity} evaluates spectral
budgets on complete Study~2 banks using the score $C_{16}/16$.

\subsection{From a fitted predictor to an action class}

We now define a common loss for comparing three classes of
linear maps on a fixed representation bank.
Throughout this comparison, we hold the observations, the
prescribed input action, and the evaluation metric fixed.
Let $r$ denote the dimension of the retained coordinate space,
and write $z_{i,t}\in\R^{1\times r}$ for the observation of
state $i$ in template $t$.
Expectations average over a fixed template distribution
$P_T$ and an independent uniform state in $S$.
For a finite bank with complete templates
$\mathcal T_B$, we use the empirical average
\[
 \mathbb E[\varphi(t,i)]
 =
 \frac{1}{|\mathcal T_B|\,n}
 \sum_{t\in\mathcal T_B}\sum_{i\in S}\varphi(t,i),
 \qquad n=|S|.
\]
The reported Study~2 class-risk calculations use the
128 complete test-template variants in each condition.
Group averages include every element, including the identity,
and retain fixed-point pairs.

Assume that the uncentered second moment is positive definite.
Define the whitened representation rows by
\[
 \Sigma=\mathbb E[z^\top z]\succ0,
 \qquad
 x_{i,t}=z_{i,t}\Sigma^{-1/2},
 \qquad
 x_g(i,t)=x_{g(i),t},
 \qquad
 \mathbb E[x^\top x]=I_r.
\]
Here $x$ denotes a whitened representation.
For the empirical class-risk calculations, $\Sigma$ is
computed on the bank being characterized.
On the original bank, this second moment includes the
contribution of the template means.
The primary analysis uses $r=q=16$.
The eight-coordinate sensitivity analysis uses its declared
truncated space and its own second moment
(\cref{app:spectral-dimension-sensitivity}).
Each retained space must satisfy the positive-definiteness
condition above.

For matrices $A_g\in\R^{r\times r}$ shared across templates,
define the common prediction loss
\begin{equation}
 \mathcal L((A_g)_g)
 =
 \frac{1}{r|G|}
 \sum_{g\in G}\mathbb E\|xA_g-x_g\|_2^2.
 \label{eq:class-loss-main}
\end{equation}
Since $\mathbb E\|x\|_2^2=r$, the factor $1/r$ normalizes
by the total representation energy in this whitened metric.
Allowing a separate unrestricted linear map for each element
gives the ordinary least-squares optimum $\eta$.
Requiring each map to be orthogonal gives
$L_{\mathrm{orth}}$, while still allowing the element
matrices to be chosen independently.
Finally, define the class of exact row actions
\[
 \mathcal R_r(G)
 =
 \left\{
 \rho:G\to O(r):
 \rho_e=I_r,\ 
 \rho_g\rho_h=\rho_{hg}
 \text{ for all }g,h\in G
 \right\},
\]
where
$O(r)=\{Q\in\R^{r\times r}:Q^\top Q=I_r\}$
includes rotations and reflections.
The reversed multiplication order follows the row-vector
convention in \cref{eq:convention}.
Transposing the matrices gives the usual column-vector
representation convention.
This class permits nonfaithful representations, including
the action that assigns $I_r$ to every group element.

The three class optima are
\[
 \begin{aligned}
 \eta
   &=\min_{(A_g)_g\in(\R^{r\times r})^G}
       \mathcal L((A_g)_g),\\
 L_{\mathrm{orth}}
   &=\min_{(Q_g)_g\in O(r)^G}
       \mathcal L((Q_g)_g),\\
 L^\star
   &=\min_{\rho\in\mathcal R_r(G)}
       \mathcal L(\rho).
 \end{aligned}
\]
Under the unit-second-moment condition, the independent
least-squares minimizers are
$A_g^\star=\mathbb E[x^\top x_g]$.
Both orthogonal classes are nonempty and compact, and the
loss is continuous, so their minima are attained as well.
Class inclusion gives
\[
 \eta\le L_{\mathrm{orth}}\le L^\star.
\]
The difference $L_{\mathrm{orth}}-\eta$ measures the cost
of orthogonality, while $L^\star-L_{\mathrm{orth}}$
measures the further cost of enforcing a common group law.
Both costs refer to the same bank, target pairs, and metric.
We compute $\eta$ and $L_{\mathrm{orth}}$ directly and
bound $L^\star$ using the spectral construction in the
next section.
We use OLS to denote ordinary least squares.

Constructed-action transfer evaluates a particular feasible
action obtained from training observations.
It applies the fixed training-derived matrices to held-out
observations using the frozen training whitening.
Its prediction loss is therefore evaluated in the
training-defined metric.
The finite-bank class-risk calculations above use the
evaluated bank's own second moment.

The registration comparison varies the action on state
labels while holding the representation bank fixed.
A \emph{registration} is a homomorphism
$\alpha:G\to\mathrm{Sym}(S)$, with
$g(i)=\alpha(g)(i)$.
For a state permutation $\pi$, the conjugate action
\[
 \alpha_\pi(g)=\pi\,\alpha(g)\,\pi^{-1}
\]
defines coherent alternative targets on the same labeled
source observations.
In the faithful, transitive $C_7$ and $D_4$ comparisons,
we enumerate distinct image subgroups $\alpha(G)$.
Reparameterizing a registration by a group automorphism
reindexes the uniform group average and preserves the
spectral budget.
For example, the one-day and two-day weekday shifts
generate the same image subgroup and receive the same
full-average spectral score.
The spectral construction below compares the energy
captured for each candidate's target correspondences.
An ordering of the exact-action optima requires separate
justification from the corresponding risk bounds.

The centering analysis in
\cref{app:template-centering-sensitivity} additionally
examines the dependence of these spectral comparisons
on the state-invariant component.
Its original and rewhitened centered observations satisfy
the unit-second-moment condition used here.
For centered observations retaining the original whitening,
$C_{16}/16$ uses the original total energy as its scale.
That treatment is reported through its spectral budgets
and registration ranks.
Table~\ref{tab:metric-contract} summarizes the prediction
targets, normalizations, and averaging rules.

\begin{table}[t]
\caption{
Evaluation targets and normalizations.
Each row specifies its own prediction object or spectral
quantity and averaging rule.
Full-state averages retain fixed-point pairs.
In the class-risk analysis, $r$ is the declared retained
dimension and the bank satisfies the stated nonsingularity
condition.
}
\label{tab:metric-contract}
\centering
\small
\begin{tabular}{
 @{}p{.19\textwidth}
 p{.40\textwidth}
 p{.33\textwidth}@{}
}
\toprule
Target & Loss or normalization & Evaluation scope \\
\midrule

{\raggedright
Structural prediction and consistency\par}
&
$D(A,B)$ divides each row's squared discrepancy by the
comparison row's squared norm plus $10^{-12}$.
&
Study~1, complete test-template state blocks.
Prediction and consistency use the element or element-pair
averages defined above.
\\[3pt]

Coordinate prediction
&
Squared error divided by
$E_z=\mathbb E_{\rm train}
 \|z-\bar z_{\rm train}\|_2^2$.
&
Study~2 fitted predictors.
Scoring uses changed-state pairs, with action weighting
specified by the experiment.
\\[3pt]

{\raggedright
Context-difference transport\par}
&
Squared error divided by
$E_u=\mathbb E_{\rm train}\|u\|_2^2$,
where $u$ is the difference between paired
same-state contexts.
&
Study~2 difference prediction.
Scoring uses changed-state pairs, with the experiment's
specified action weighting.
\\[3pt]

Action-class risks
&
$\mathcal L$ in \cref{eq:class-loss-main},
using the evaluated bank's second-moment whitening
and division by $r$.
&
$\eta$, $L_{\mathrm{orth}}$, and bounds on $L^\star$
on one bank.
Uniform averages over complete state blocks and all
group elements.
\\[3pt]

{\raggedright
Constructed-action transfer\par}
&
Squared prediction loss divided by $r$,
using the fixed training whitening and constructed matrices.
&
A particular feasible action evaluated on held-out
observations.
Complete state blocks and all group elements.
\\[3pt]

{\raggedright
Spectral budget and registration\par}
&
$C_r(K)/r$, equal to $\gamma_r$ under the
unit-second-moment condition.
The fixed-whitening centered treatment uses the original
energy scale.
&
Finite-bank spectral compatibility and candidate-image
rankings.
The centering treatments are specified in
\cref{app:template-centering-sensitivity}.
\\

\bottomrule
\end{tabular}
\end{table}

% Keep the evaluation-contract table before the theory section.
\FloatBarrier
